\section{Introduction}

\subsection{Objet du document}

Ce document décrit comment l'application \emph{Tado pour Pebble} est conçue pour satisfaire
le cahier des charges TPB-SyRS [P1]. Le cahier des charges dit \emph{ce que} l'application
fait~; le présent document dit \emph{comment} elle est découpée, quelles briques portent
quelle exigence, comment elles échangent, et pourquoi ces choix ont été faits plutôt que
d'autres.

Il s'adresse à celui qui écrira le code (le découpage, les protocoles, les règles de calcul)
et à celui qui le relira (les décisions, leurs raisons, leurs conséquences). Il n'est pas le
code~: il s'arrête là où l'écriture devient évidente.

\subsection{Périmètre}

\paragraph{Dans le périmètre.}
\begin{itemize}
  \item le découpage en programmes (montre, téléphone) et en composants~;
  \item les échanges avec Tado, leur coût en requêtes et la politique qui les contient~;
  \item le protocole entre la montre et le téléphone~;
  \item les modèles de données manipulés à chaque étape~;
  \item l'adaptation à chacune des six plateformes Pebble~;
  \item les comportements dynamiques (ouverture, commande, liaison, erreurs)~;
  \item la traçabilité entre exigences et composants, et les retours que cette conception
        fait au cahier des charges (§\ref{sec:retours}).
\end{itemize}

\paragraph{Hors périmètre.} Le code source, les procédures de publication sur le catalogue
de montres, et tout ce qui relève des besoins fonctionnels, déjà spécifiés.

\subsection{Principes de rédaction}

\paragraph{Décisions.} Chaque choix structurant est une \emph{décision}, identifiée
\texttt{DEC-nn}, présentée dans un cartouche (contexte, décision, alternatives écartées,
conséquences, exigences concernées). Une décision n'est jamais supprimée~: si elle change,
on la déclare remplacée.

\paragraph{Composants.} Les composants du téléphone sont identifiés \texttt{J-nn}, ceux de la
montre \texttt{M-nn}, les messages \texttt{MSG-nn}. Les exigences du cahier des charges sont
citées par leur identifiant \texttt{TPB-\textit{DOM}-\textit{nnn}}.

\paragraph{Formes verbales.} Les mêmes que dans le cahier des charges~: \enquote{doit}
exprime une obligation, \enquote{devrait} une recommandation, l'indicatif un fait de
conception.

\paragraph{Points ouverts.} Les questions qui restent à trancher sont signalées dans le corps
du document et récapitulées en annexe (§\ref{sec:points-ouverts}).

\subsection{Définitions complémentaires}

\begin{longtable}{@{}>{\bfseries}p{42mm}p{118mm}@{}}
  Watchapp & Programme qui s'exécute sur la montre, écrit en C avec le SDK Pebble. \\
  PebbleKit JS (pkjs) & Programme JavaScript qui s'exécute sur le téléphone, dans l'application Pebble qui accompagne la montre. Il peut faire des requêtes réseau et converse avec la watchapp. \\
  AppMessage & Mécanisme de messages entre la watchapp et le pkjs, par paires clé-valeur dans des tampons de taille fixe. \\
  Vue (modèle de vue) & Données déjà préparées pour un écran~: noms abrégés, températures arrondies, niveaux calculés. La montre n'affiche que des vues. \\
  TTL & Durée pendant laquelle une donnée mise en cache est tenue pour assez fraîche pour être affichée sans nouvelle requête. \\
  Consigne & Température demandée à une pièce. \\
  Forçage (overlay) & Ordre manuel qui remplace la programmation d'une pièce, avec une condition de fin. \\
  Zone & Terme de l'API Tado pour une pièce. \\
\end{longtable}
